\subsubsection{Subproceso 1: Acreditación y Pre-requisitos}

\begin{itemize}
    \item \textbf{Regla de pago obligatorio del Autoseguro Estudiantil:} Para habilitar su matrícula regular, el estudiante debe haber pagado antes el aporte al Fondo de Autoseguro Estudiantil del ciclo correspondiente~\cite[Nota~1]{tupa2024}, \cite[nums.~6.1.1 y 6.2]{autoseguro2023}.

    \item \textbf{Regla de no tener adeudos en la Facultad:} El alumno no debe presentar deudas pendientes de libros, materiales de enseñanza o equipos en su Facultad~\cite[Nota~1]{tupa2024}, \cite[Arts.~46\textdegree{} y 47\textdegree{}]{reglamento2022}. Por el contrario, no haber pagado el carné universitario o tener multas electorales no impide matricularse~\cite[Art.~47\textdegree{}, incs.~h e i]{reglamento2022}.

    \item \textbf{Regla de cierre previo de notas:} Las actas de evaluación del semestre anterior deben estar cerradas y sin reclamos de notas pendientes al menos tres (03) días útiles antes de iniciar la matrícula regular~\cite[Art.~6\textdegree{}]{reglamento2022}.

    \item \textbf{Regla de turno por orden de mérito:} El grupo y horario de entrada a la matrícula regular se asignan según el promedio ponderado de los dos últimos semestres en que el alumno tuvo matrícula regular, ordenando de mayor a menor promedio~\cite[Art.~9\textdegree{}, inc.~a; Art.~11\textdegree{}, inc.~m]{reglamento2022}.

    \item \textbf{Regla de turno preferencial para promoción (egresantes):} Los estudiantes a los que les falten 30 créditos o menos para terminar su carrera tienen prioridad en un grupo preferencial antes de los turnos regulares, siempre que lo soliciten a su Escuela Profesional hasta tres (03) días hábiles antes de la matrícula~\cite[Art.~42\textdegree{}]{reglamento2022}.

    \item \textbf{Regla de anticipación para cambios de plan y convalidaciones:} Si un estudiante con código activo cambia de plan de estudios o pide exonerar una electiva desaprobada, el trámite debe solicitarse hasta diez (10) días hábiles antes de la semana de matrícula, y la DIRCE debe procesar las convalidaciones antes de abrir los turnos para que la ficha refleje los cursos correctos~\cite[Arts.~17\textdegree{} y 43\textdegree{}]{reglamento2022}.

    \item \textbf{Regla de tutoría previa para activar el código en riesgo académico:} El estudiante con asignaturas desaprobadas dos o tres veces debe pasar por tutoría obligatoria antes de la matrícula para que su tutor autorice la activación de su código~\cite[Arts.~8\textdegree{} y 48\textdegree{}]{reglamento2022}.
\end{itemize}
